\documentclass[10pt,a4paper]{amsart}
\usepackage{amsmath,amssymb,amsthm,graphicx}
% Retained minimal preamble for the Stacks Project chapter sites-cohomology.tex (Cohomology on Sites).
% Verbatim excerpt of the chapter preamble: Xy-pic setup (native preamble lines 14-18) and the native
% theorem declarations with the mathematical macros (native preamble lines 154-217). Copied mechanically;
% nothing added, removed or redefined. The wrapper class and the amsmath/amssymb/amsthm/graphicx base
% packages are supplied by the assembler, so the class/externaldocument part of the native preamble is omitted.
% For commutative diagrams we use Xy-pic
\usepackage[all]{xy}

% We use 2cell for 2-commutative diagrams.
\xyoption{2cell}
% Theorem environments.
%
\theoremstyle{plain}
\newtheorem{theorem}[subsection]{Theorem}
\newtheorem{proposition}[subsection]{Proposition}
\newtheorem{lemma}[subsection]{Lemma}

\theoremstyle{definition}
\newtheorem{definition}[subsection]{Definition}
\newtheorem{example}[subsection]{Example}
\newtheorem{exercise}[subsection]{Exercise}
\newtheorem{situation}[subsection]{Situation}

\theoremstyle{remark}
\newtheorem{remark}[subsection]{Remark}
\newtheorem{remarks}[subsection]{Remarks}

\numberwithin{equation}{subsection}

% Macros
%
\def\lim{\mathop{\mathrm{lim}}\nolimits}
\def\colim{\mathop{\mathrm{colim}}\nolimits}
\def\Spec{\mathop{\mathrm{Spec}}}
\def\Hom{\mathop{\mathrm{Hom}}\nolimits}
\def\Ext{\mathop{\mathrm{Ext}}\nolimits}
\def\SheafHom{\mathop{\mathcal{H}\!\mathit{om}}\nolimits}
\def\SheafExt{\mathop{\mathcal{E}\!\mathit{xt}}\nolimits}
\def\Sch{\mathit{Sch}}
\def\Mor{\mathop{\mathrm{Mor}}\nolimits}
\def\Ob{\mathop{\mathrm{Ob}}\nolimits}
\def\Sh{\mathop{\mathit{Sh}}\nolimits}
\def\NL{\mathop{N\!L}\nolimits}
\def\CH{\mathop{\mathrm{CH}}\nolimits}
\def\proetale{{pro\text{-}\acute{e}tale}}
\def\etale{{\acute{e}tale}}
\def\QCoh{\mathit{QCoh}}
\def\Ker{\mathop{\mathrm{Ker}}}
\def\Im{\mathop{\mathrm{Im}}}
\def\Coker{\mathop{\mathrm{Coker}}}
\def\Coim{\mathop{\mathrm{Coim}}}

% Boxtimes
%
\DeclareMathSymbol{\boxtimes}{\mathbin}{AMSa}{"02}

%
% Macros for moduli stacks/spaces
%
\def\QCohstack{\mathcal{QC}\!\mathit{oh}}
\def\Cohstack{\mathcal{C}\!\mathit{oh}}
\def\Spacesstack{\mathcal{S}\!\mathit{paces}}
\def\Quotfunctor{\mathrm{Quot}}
\def\Hilbfunctor{\mathrm{Hilb}}
\def\Curvesstack{\mathcal{C}\!\mathit{urves}}
\def\Polarizedstack{\mathcal{P}\!\mathit{olarized}}
\def\Complexesstack{\mathcal{C}\!\mathit{omplexes}}
% \Pic is the operator that assigns to X its picard group, usage \Pic(X)
% \Picardstack_{X/B} denotes the Picard stack of X over B
% \Picardfunctor_{X/B} denotes the Picard functor of X over B
\def\Pic{\mathop{\mathrm{Pic}}\nolimits}
\def\Picardstack{\mathcal{P}\!\mathit{ic}}
\def\Picardfunctor{\mathrm{Pic}}
\def\Deformationcategory{\mathcal{D}\!\mathit{ef}}
\usepackage{cleveref}
\crefname{originalequation}{original equation}{original equations}
\Crefname{originalequation}{Original equation}{Original equations}
\makeatletter
\newlabel{equation-eval}{{08I2}{source}{}{sourcepointer-1}{}}
\newlabel{equation-eval@cref}{{[originalequation][0][]08I2}{[1][0][]source}{}{}{}}
\makeatother
\crefname{originallemma}{original lemma}{original lemmas}
\Crefname{originallemma}{Original lemma}{Original lemmas}
\makeatletter
\newlabel{lemma-Rhom-strictly-perfect}{{08DM}{source}{}{sourcepointer-2}{}}
\newlabel{lemma-Rhom-strictly-perfect@cref}{{[originallemma][0][]08DM}{[1][0][]source}{}{}{}}
\makeatother
\crefname{originallemma}{original lemma}{original lemmas}
\Crefname{originallemma}{Original lemma}{Original lemmas}
\makeatletter
\newlabel{lemma-dual}{{08I1}{source}{}{sourcepointer-3}{}}
\newlabel{lemma-dual@cref}{{[originallemma][0][]08I1}{[1][0][]source}{}{}{}}
\makeatother
\crefname{originallemma}{original lemma}{original lemmas}
\Crefname{originallemma}{Original lemma}{Original lemmas}
\makeatletter
\newlabel{lemma-internal-hom-evaluate}{{0A67}{source}{}{sourcepointer-4}{}}
\newlabel{lemma-internal-hom-evaluate@cref}{{[originallemma][0][]0A67}{[1][0][]source}{}{}{}}
\makeatother
\crefname{originallemma}{original lemma}{original lemmas}
\Crefname{originallemma}{Original lemma}{Original lemmas}
\makeatletter
\newlabel{lemma-restriction-RHom-to-U}{{08DL}{source}{}{sourcepointer-5}{}}
\newlabel{lemma-restriction-RHom-to-U@cref}{{[originallemma][0][]08DL}{[1][0][]source}{}{}{}}
\makeatother
\title{The derived dual of a perfect complex on a ringed site is perfect and involutive, with functorial identifications against internal hom and Hom in the derived category}
\author{The Stacks Project Authors}
\date{stacks-project:sites-cohomology.tex:local-snapshot-20261011; GNU FDL 1.2 or later}
\begin{document}
\maketitle
\noindent\textit{Editorial note.} All delivered mathematical text is the author's own, copied from the pinned Stacks Project chapter file sites-cohomology.tex; no statement or proof step was rewritten, reordered or removed.\par\medskip
\noindent\textit{Editorial note.} The counted proof is the author's complete argument: the proof is the author's complete argument: it reduces the perfection statement, the involution and the two identifications to a local computation, and performs that computation on a strictly perfect representative using the internal-hom lemma for strictly perfect complexes.\par\medskip
\noindent\textit{Editorial note.} The label lemma-dual-perfect-complex names the counted statement, and the proof environment immediately following it in the source is delivered verbatim as its proof.\par\medskip
\noindent\textit{Editorial note.} The statement's two displayed isomorphisms are delivered verbatim, including the internal hom and Hom notation of the chapter.\par\medskip
\noindent\textit{Editorial note.} The preamble is the chapter's own retained minimal preamble (sources/stacks-sites-cohomology/retained-preamble.tex), a verbatim excerpt of the native preamble: the Xy-pic setup and the native theorem declarations with the mathematical macros.\par\medskip
\noindent\textit{Allowed standard prerequisites (editorial).}\begin{itemize}
\item Perfect objects of D(O) on a ringed site. Same chapter, the definition of a perfect object as one that is locally represented by a strictly perfect complex. Supplies the locality reduction: perfection of K is checked after restricting along a covering of the site.
\item Derived dual K\textasciicircum{}vee = RSheafHom(K, O). Same chapter, the internal hom into the structure sheaf and the evaluation map equation-eval. Supplies the object K\textasciicircum{}vee whose perfection and involution are asserted, and the evaluation map used in the identification.
\item Internal hom commutes with restriction to C/U. Same chapter, the restriction compatibility of RSheafHom used without further mention at the start of the proof. Supplies the passage from the global assertion to the local one over the members of a covering.
\end{itemize}\medskip
\ifdefined\hypertarget\hypertarget{sourcepointer-1}{}\fi
\noindent\textit{Source reference.} Original equation 08I2: sources/stacks-sites-cohomology/pinned-original-source.tex, source lines 9119--9119. Prerequisite pointer invoked at proof line(s) 13082; it names a standard fact of the same chapter that the counted argument uses, and every step of the counted proof built on it is retained verbatim in the delivered body.\par\medskip
\ifdefined\hypertarget\hypertarget{sourcepointer-2}{}\fi
\noindent\textit{Source reference.} Original lemma 08DM: sources/stacks-sites-cohomology/pinned-original-source.tex, source lines 11912--11912. Prerequisite pointer invoked at proof line(s) 13095; it names a standard fact of the same chapter that the counted argument uses, and every step of the counted proof built on it is retained verbatim in the delivered body.\par\medskip
\ifdefined\hypertarget\hypertarget{sourcepointer-3}{}\fi
\noindent\textit{Source reference.} Original lemma 08I1: sources/stacks-sites-cohomology/pinned-original-source.tex, source lines 9114--9114. Prerequisite pointer invoked at proof line(s) 13081; it names a standard fact of the same chapter that the counted argument uses, and every step of the counted proof built on it is retained verbatim in the delivered body.\par\medskip
\ifdefined\hypertarget\hypertarget{sourcepointer-4}{}\fi
\noindent\textit{Source reference.} Original lemma 0A67: sources/stacks-sites-cohomology/pinned-original-source.tex, source lines 8965--8965. Prerequisite pointer invoked at proof line(s) 13078; it names a standard fact of the same chapter that the counted argument uses, and every step of the counted proof built on it is retained verbatim in the delivered body.\par\medskip
\ifdefined\hypertarget\hypertarget{sourcepointer-5}{}\fi
\noindent\textit{Source reference.} Original lemma 08DL: sources/stacks-sites-cohomology/pinned-original-source.tex, source lines 8934--8934. Prerequisite pointer invoked at proof line(s) 13067; it names a standard fact of the same chapter that the counted argument uses, and every step of the counted proof built on it is retained verbatim in the delivered body.\par\medskip
\begin{lemma}
\label{lemma-dual-perfect-complex}
Let $(\mathcal{C}, \mathcal{O})$ be a ringed site.
Let $K$ be a perfect object of $D(\mathcal{O})$.
Then $K^\vee = R\SheafHom(K, \mathcal{O})$ is a
perfect object too and $(K^\vee)^\vee \cong K$. There are
functorial isomorphisms
$$
M \otimes^\mathbf{L}_\mathcal{O} K^\vee = R\SheafHom_\mathcal{O}(K, M)
$$
and
$$
H^0(\mathcal{C}, M \otimes^\mathbf{L}_\mathcal{O} K^\vee) =
\Hom_{D(\mathcal{O})}(K, M)
$$
for $M$ in $D(\mathcal{O})$.
\end{lemma}

\begin{proof}
We will us without further mention that formation of internal hom commutes
with restriction (Lemma \ref{lemma-restriction-RHom-to-U}). Let $U$
be an arbitrary object of $\mathcal{C}$. To check that
$K^\vee$ is perfect, it suffices to show that there exists a covering
$\{U_i \to U\}$ such that $K^\vee|_{U_i}$ is perfect for all $i$.
There is a canonical map
$$
K = R\SheafHom(\mathcal{O}_X, \mathcal{O}_X)
\otimes_{\mathcal{O}_X}^\mathbf{L} K \longrightarrow
R\SheafHom(R\SheafHom(K, \mathcal{O}_X), \mathcal{O}_X) =
(K^\vee)^\vee
$$
see Lemma \ref{lemma-internal-hom-evaluate}. It suffices to prove there
is a covering $\{U_i \to U\}$ such that the restriction of this map
to $\mathcal{C}/U_i$ is an isomorphism for all $i$.
By Lemma \ref{lemma-dual} to see the final statement it suffices to check
that the map (\ref{equation-eval})
$$
M \otimes^\mathbf{L}_\mathcal{O} K^\vee
\longrightarrow
R\SheafHom(K, M)
$$
is an isomorphism. This is a local question as well (in the sense above).
Hence it suffices to prove the lemma when $K$ is represented
by a strictly perfect complex.

\medskip\noindent
Assume $K$ is represented by the strictly perfect complex
$\mathcal{E}^\bullet$. Then it follows from
Lemma \ref{lemma-Rhom-strictly-perfect}
that $K^\vee$ is represented by the complex whose terms are
$(\mathcal{E}^n)^\vee =
\SheafHom_\mathcal{O}(\mathcal{E}^n, \mathcal{O})$
in degree $-n$. Since $\mathcal{E}^n$ is a direct summand of a finite
free $\mathcal{O}$-module, so is $(\mathcal{E}^n)^\vee$.
Hence $K^\vee$ is represented by a strictly perfect complex too 
and we see that $K^\vee$ is perfect.
The map $K \to (K^\vee)^\vee$ is an isomorphism as it is given up
to sign by the evaluation maps
$\mathcal{E}^n \to ((\mathcal{E}^n)^\vee)^\vee$ which are
isomorphisms. To see that (\ref{equation-eval}) is an isomorphism, represent
$M$ by a K-flat complex $\mathcal{F}^\bullet$.
By Lemma \ref{lemma-Rhom-strictly-perfect} the complex
$R\SheafHom(K, M)$ is represented by the complex with terms
$$
\bigoplus\nolimits_{n = p + q}
\SheafHom_\mathcal{O}(\mathcal{E}^{-q}, \mathcal{F}^p)
$$
On the other hand, the object $M \otimes^\mathbf{L}_\mathcal{O} K^\vee$
is represented by the complex with terms
$$
\bigoplus\nolimits_{n = p + q}
\mathcal{F}^p \otimes_\mathcal{O} (\mathcal{E}^{-q})^\vee
$$
Thus the assertion that (\ref{equation-eval}) is an isomorphism
reduces to the assertion that the canonical map
$$
\mathcal{F}
\otimes_\mathcal{O}
\SheafHom_\mathcal{O}(\mathcal{E}, \mathcal{O})
\longrightarrow
\SheafHom_\mathcal{O}(\mathcal{E}, \mathcal{F})
$$
is an isomorphism when $\mathcal{E}$ is a direct summand of a finite
free $\mathcal{O}$-module and $\mathcal{F}$ is any $\mathcal{O}$-module.
This follows immediately from the corresponding statement when
$\mathcal{E}$ is finite free.
\end{proof}

\end{document}
